\documentclass[12pt,a4paper]{article}

% --------------------------------------------------
% Packages
% --------------------------------------------------
\usepackage[utf8]{inputenc}
\usepackage{graphicx}
\usepackage{amsmath}
\usepackage{hyperref}
\usepackage{geometry}
\usepackage{booktabs}
\usepackage{caption}
\usepackage{subcaption}
\usepackage{float}
\usepackage{listings}
\usepackage{xcolor}

\geometry{margin=1in}

% Code style
\lstset{
  basicstyle=\ttfamily\footnotesize,
  breaklines=true,
  frame=single,
  backgroundcolor=\color{gray!10},
  keywordstyle=\color{blue},
  commentstyle=\color{green!50!black},
  stringstyle=\color{orange}
}

% -----------------------------------------------------------------------------
% Title
% -----------------------------------------------------------------------------
\title{Hydropower Site Suitability}
\author{Alin-Ioan Alexandru \\
        Stefan-Cristian Grecu \\}
\date{\today}

% -----------------------------------------------------------------------------
% Document
% -----------------------------------------------------------------------------
\begin{document}
\maketitle

\tableofcontents
\newpage

% -----------------------------------------------------------------------------
\section{Introduction}

This document describes the methodology, data processing pipeline, model
training, and evaluation for the project focused on predicting site suitability
for new hydropower plants using satellite imagery and deep learning segmentation
models.
\\
\\
The objective is to identify optimal locations for new hydropower plants by
analyzing a variety of environmental factors using supervised learning. While
there is limited available research that directly uses machine learning models
in the analysis of site suitability for hydropower sites, our method, with the
use of satellite imagery, has the advantage that it can potentially uncovere
suitable locations that may not be considered by typical manual or rule-based
analyses.
\\
\\
We decided to use two approaches, resulting in two ML pipelines. The two
pipelines start the same by making requests to the Sentinel Hub API to acquire
the satellite data. The first approach models the gathered data in a tabular
format and uses traditional ML algorithms (Random Forest, XGBoost) to make
predictions. The second approach models the data as images with multiple
channels and uses a deep learning segmentation model (U-Net) to make predictions.

% -----------------------------------------------------------------------------
\section{Data Sources and Gathering}

The resulting dataset consists of about 3000 positive samples consisting of
locations of hydropower plants from Europe and 5000 negative samples of random
locations from European rivers.

\subsection{Satellite Data}

Using the Sentinel Hub API, we gather Sentinel-2 satellite imagery for the
regions of interest. We focus on bands B3, B4, B8 and B11, which cover the
visible and near-infrared spectrum.
\\
\\
The following spectral indices are calculated:
\begin{itemize}
  \item NDVI (Normalized Difference Vegetation Index) - used to assess
  vegetation health and density. Vegetation cover can be used to assess land
  suitability for hydropower plants.
  \item NDWI (Normalized Difference Water Index) - used to identify water
  bodies and assess their extent. Water detection is important in identifying
  potential sites for hydropower plants.
  \item NDBI (Normalized Difference Built-up Index) - used to detect built-up
  areas and urbanization. This helps in avoiding locations that are too close
  to urban areas.
  \item MNDWI (Modified Normalized Difference Water Index) - used to enhance
  water features in urban areas.
\end{itemize}

\subsection{Topographic Data}

Topographic data is obtained from the Copernicus Digital Elevation Model (DEM).
The DEM provides elevation data which is crucial for assessing the potential
energy generation capacity of a site and the potential of forming and
accumulation lakes near the plants. From the DEM, we also derive slope data,
which helps in understanding the terrain and its suitability for hydropower
infrastructure.

\subsection{Hydrological and Precipitation Data}

Hydrological data, specifically river discharge data, is obtained from the
Hydrorives database and precipitation data is gathered using the Open-Meteo
API. River discharge is a key factor in determining the potential energy
output of a hydropower plant, while precipitation data helps in understanding
the water availability in the region.

\subsection{Labels}
Positive samples: near existing hydropower plants. \\
Negative samples: random locations at least 30 kilometers away from any
hydropower plant. \\
Mask values (for the segmentation) and the classes used are as following:\\
\texttt{1 = suitable}, \texttt{0 = unsuitable}.

% -----------------------------------------------------------------------------
\section{Data Preprocessing}

\subsection{Preprocessing for Tabular Data}
Preprocessing steps:
\begin{itemize}
  \item Splitting data: Performed train-test split (e.g., 80-20).
  \item Handling missing values: Imputed missing values using mean strategy.
  \item Outlier detection: Identified and treated outliers using IQR method and
  replaced them using the same imputation strategy as for the missing values.
  \item Normalization: Scaled features to a standard range.
\end{itemize}

\subsection{Preprocessing for Image Data}
Feature cube construction: shape $(H, W, C)$ with channels for spectral indices,
DEM, slope, discharge, precipitation.

% -----------------------------------------------------------------------------
\section{Machine Learning Models}

\subsection{Tabular Data Models}

For the tabular representation of the dataset, we trained and compared three
widely used ensemble learning algorithms:

\begin{itemize}
  \item \textbf{Random Forest}: An ensemble of decision trees using bootstrap
  sampling and feature bagging. It provides robust performance, reduces
  overfitting, and allows for feature importance ranking.
  \item \textbf{MLP (Multi-Layer Perceptron)}: A feedforward neural network with
  multiple layers and non-linear activation functions. It can model complex
  relationships in the data and is suitable for tabular data.
  \item \textbf{XGBoost}: An optimized implementation of gradient boosting
  decision trees. It generally provides higher accuracy, faster convergence,
  and better handling of imbalanced data.
\end{itemize}

\subsection{Segmentation Model}

For the image-based approach, we used a \textbf{U-Net} architecture implemented
via the Segmentation Models PyTorch (SMP) library. The encoder was initialized
with weights pre-trained on ImageNet to accelerate convergence.

\begin{itemize}
  \item \textbf{Input}: Multi-channel feature cubes composed of spectral
  indices, DEM, slope, discharge, and precipitation layers.
  \item \textbf{Loss function}: Combination of Cross-Entropy Loss and Dice Loss
  to balance pixel-wise classification and overlap quality.
  \item \textbf{Optimizer}: Adam optimizer with an initial learning rate of
  $1 \times 10^{-3}$ and learning rate scheduling.
  \item \textbf{Training}: Conducted for 20--30 epochs with a batch size of 16.
\end{itemize}

The model outputs a binary mask where each pixel is classified as either
suitable or unsuitable for hydropower site placement.

% -----------------------------------------------------------------------------
\section{Evaluation and Results}

\subsection{Tabular Data Models}

\subsection{segmentation Model}

% -----------------------------------------------------------------------------
\section{Conclusion}

% -----------------------------------------------------------------------------

\end{document}
